\documentclass[11pt,a4paper]{article}
\usepackage[utf8]{vietnam}
\usepackage{amsmath,amssymb}
\usepackage{tabularx}
\usepackage{tikz}
\usepackage{xcolor}
\usepackage{fancyhdr}
\usepackage[top=1.8cm, bottom=2.0cm, left=1.8cm, right=1.8cm]{geometry}

\setlength{\parskip}{3pt}
\setlength{\parindent}{0pt}

\pagestyle{fancy}
\fancyhf{}
\renewcommand{\headrulewidth}{0.4pt}
\lhead{\footnotesize \textbf{CHÙA HƯƠNG SƠN -- HƯỚNG DẪN CHẤM THI}}
\rhead{\footnotesize \textbf{ĐÁP ÁN CHI TIẾT ĐỀ THI THỬ SỐ 03}}
\cfoot{\thepage}

\begin{document}

\begin{center}
{\large\bfseries ĐÁP ÁN \& HƯỚNG DẪN CHẤM ĐỀ THI THỬ CUỐI HỌC KỲ I (ĐỀ SỐ 03)}\\[1.5mm]
{\bfseries ĐƠN VỊ: CHÙA HƯƠNG SƠN \quad|\quad MÔN: TOÁN -- LỚP 9 \quad|\quad NĂM HỌC 2026 -- 2027}\\[1mm]
\textit{(Chuẩn cấu trúc định dạng mới của Sở Giáo dục và Đào tạo)}
\end{center}

\vspace{1mm}
\hrule
\vspace{3mm}

\noindent
\textbf{\large PHẦN I. TRẮC NGHIỆM KHÁCH QUAN (7,0 ĐIỂM)}

\vspace{2mm}
\noindent
\textbf{1. Bảng đáp án Phần 1: Trắc nghiệm nhiều lựa chọn (3,0 điểm)} \textit{(Mỗi câu đúng được 0,25 điểm)}
\begin{center}
\begin{tabular}{|c|c|c|c|c|c|c|c|c|c|c|c|c|}
\hline
\textbf{Câu} & \textbf{1} & \textbf{2} & \textbf{3} & \textbf{4} & \textbf{5} & \textbf{6} & \textbf{7} & \textbf{8} & \textbf{9} & \textbf{10} & \textbf{11} & \textbf{12} \\
\hline
\textbf{Đáp án} & \textbf{A} & \textbf{A} & \textbf{A} & \textbf{A} & \textbf{B} & \textbf{A} & \textbf{A} & \textbf{A} & \textbf{A} & \textbf{A} & \textbf{A} & \textbf{A} \\
\hline
\end{tabular}
\end{center}

\vspace{2mm}
\noindent
\textbf{2. Bảng đáp án Phần 2: Trắc nghiệm Đúng -- Sai (2,0 điểm)}
\begin{center}
\begin{tabular}{|c|c|c|c|c|l|}
\hline
\textbf{Câu} & \textbf{Ý a} & \textbf{Ý b} & \textbf{Ý c} & \textbf{Ý d} & \textbf{Ghi chú đáp án chuẩn} \\
\hline
\textbf{Câu 13} & \textbf{Đ} & \textbf{S} & \textbf{Đ} & \textbf{Đ} & Ý b) sai vì vế phải phải bằng $\dfrac{1}{2}$ khối lượng công việc. \\
\hline
\textbf{Câu 14} & \textbf{Đ} & \textbf{Đ} & \textbf{Đ} & \textbf{S} & Ý d) sai vì chu vi $\Delta ACD = 2AC + CD = 2(4\sqrt{5}) + 16 \ne 48\text{ cm}$. \\
\hline
\end{tabular}
\end{center}

\vspace{2mm}
\noindent
\textbf{3. Bảng đáp án Phần 3: Trắc nghiệm trả lời ngắn (2,0 điểm)} \textit{(Mỗi câu đúng được 0,5 điểm)}
\begin{center}
\begin{tabular}{|c|l|c|c|}
\hline
\textbf{Câu} & \textbf{Nội dung câu hỏi} & \textbf{Kết quả bài làm} & \textbf{Điểm} \\
\hline
\textbf{Câu 15} & Nghiệm bất phương trình $4(x - 2) + 3 > 5(x + 1) - 8$ & \textbf{$x < 0$} & 0,5đ \\
\hline
\textbf{Câu 16} & Giá trị biểu thức $P = \sqrt{(\sqrt{6}-3)^2} + \sqrt{6}$ & \textbf{$3$} & 0,5đ \\
\hline
\textbf{Câu 17} & Chiều rộng khúc sông ($40 \cdot \tan 36^\circ$) & \textbf{$29{,}06\text{ m}$} & 0,5đ \\
\hline
\textbf{Câu 18} & Diện tích hình quạt tròn ($\dfrac{3{,}14 \times 15^2 \times 150}{360}$) & \textbf{$294{,}4\text{ cm}^2$} & 0,5đ \\
\hline
\end{tabular}
\end{center}

\vspace{3mm}
\noindent
\textbf{\large PHẦN II. TỰ LUẬN (3,0 ĐIỂM)}

\vspace{2mm}
\textbf{Câu 19 (1,0 điểm).}
\begin{itemize}
  \item[a)] \textit{(0,5 điểm)} $\begin{cases} 4x + 3y = 11 \\ 3x - y = 5 \end{cases} \Leftrightarrow \begin{cases} 4x + 3y = 11 \\ 9x - 3y = 15 \end{cases} \Leftrightarrow \begin{cases} 13x = 26 \\ y = 3x - 5 \end{cases} \Leftrightarrow \begin{cases} x = 2 \\ y = 1 \end{cases}$. \hfill \textbf{(0,5đ)}
  \item[b)] \textit{(0,5 điểm)} Với $x > 0, x \ne 1$:
  $P = \left( \dfrac{\sqrt{x}}{\sqrt{x}+1} - \dfrac{1}{\sqrt{x}(\sqrt{x}+1)} \right) : \dfrac{\sqrt{x}-1}{x} = \dfrac{x - 1}{\sqrt{x}(\sqrt{x}+1)} \cdot \dfrac{x}{\sqrt{x}-1} = \dfrac{(\sqrt{x}-1)(\sqrt{x}+1)}{\sqrt{x}(\sqrt{x}+1)} \cdot \dfrac{x}{\sqrt{x}-1} = \sqrt{x}$. \hfill \textbf{(0,5đ)}
\end{itemize}

\vspace{1mm}
\textbf{Câu 20 (1,5 điểm).}
\begin{itemize}
  \item[a)] \textit{(0,75 điểm)} Vì $MA \perp OA, MB \perp OB \Rightarrow \widehat{MAO} = \widehat{MBO} = 90^\circ$. Gọi $I$ là trung điểm $MO \Rightarrow IA = IB = IM = IO = \dfrac{MO}{2} \Rightarrow M, A, O, B$ cùng thuộc đường tròn đường kính $MO$. Lại có $MA = MB$ và $OA = OB = R \Rightarrow MO$ là đường trung trực của $AB \Rightarrow MO \perp AB$. \hfill \textbf{(0,75đ)}
  \item[b)] \textit{(0,75 điểm)} Xét $\Delta MAC$ và $\Delta MDA$ có: góc $\widehat{M}$ chung; $\widehat{MAC} = \widehat{MDA}$ (góc tạo bởi tiếp tuyến và dây cung và góc nội tiếp cùng chắn cung $AC$) $\Rightarrow \Delta MAC \backsim \Delta MDA$ (g.g) $\Rightarrow \dfrac{MA}{MD} = \dfrac{MC}{MA} \Rightarrow MA^2 = MC \cdot MD$ (đpcm). \hfill \textbf{(0,75đ)}
\end{itemize}

\vspace{1mm}
\textbf{Câu 21 (0,5 điểm).}
\begin{itemize}
  \item[a)] \textit{(0,25 điểm)} Gọi cạnh vuông góc với bờ tường là $x$ ($0 < x < 20$, m) $\Rightarrow$ cạnh song song với bờ tường là $40 - 2x$ (m). Diện tích khu đất: $S = x(40 - 2x) = -2x^2 + 40x = -2(x - 10)^2 + 200 \le 200\text{ m}^2$. Dấu ``='' xảy ra khi $x = 10\text{ m}$. Cạnh song song bờ tường là $40 - 2(10) = 20\text{ m}$. Vậy diện tích lớn nhất là $200\text{ m}^2$ khi chiều rộng $10\text{ m}$, chiều dài $20\text{ m}$. \hfill \textbf{(0,25đ)}
  \item[b)] \textit{(0,25 điểm)} Áp dụng BĐT AM-GM cho hai số dương: $a + \dfrac{1}{a} \ge 2\sqrt{a \cdot \dfrac{1}{a}} = 2$. Tương tự: $b + \dfrac{1}{b} \ge 2$ và $c + \dfrac{1}{c} \ge 2$. Cộng ba bất đẳng thức:
  \[
  Q = \left(a + \dfrac{1}{a}\right) + \left(b + \dfrac{1}{b}\right) + \left(c + \dfrac{1}{c}\right) \ge 2 + 2 + 2 = 6.
  \]
  Dấu ``='' xảy ra khi $a = b = c = 1$. Khi đó $a + b + c = 3 \le 3$ (thỏa mãn). Vậy $\min Q = 6$. \hfill \textbf{(0,25đ)}
\end{itemize}

\end{document}
